%% notes-tex/025-s3-closure/sections/6-consequences.tex

\section{What this means for training on S3\secstatus{todo}}
\label{sec:consequences}

\notesec{The identity cannot be the mechanism}
The S3 arm supervises fitness and interaction jointly on the hypothesis, stated as
untested in the training plan, that a representation shaped by fitness helps the
interaction head. If the mechanism were the model reading Eq.~\ref{eq:tau} off its
fitness predictions, its ceiling on this build would be the $r = 0.23$ that the
equation itself reaches on the stored values, below the $r = 0.40$ of the additive
ridge that sees no fitness at all. Whatever the S3 seeds show at readout, they are
not showing that; a gain would have to come from the representation, and a null would
not refute the identity, which is exact in every source.

\notesec{Per-order fitness is now logged, and per-order interaction was always logged}
The trainer reports fitness for singles, doubles and triples and interaction for
doubles and triples separately, so the readout can ask whether the model fits the
essentiality-tainted singles, which hold a value no strain has, and whether the
doubles' fitness is learned better than their interaction, which this document
predicts because the double fitness is a measurement and the double interaction is a
cross-screen difference.

\notesec{Every p-value in a merged record is unlabeled}
Nothing in the training arm reads p-values, so no run is affected. Any analysis that
does, a confidence-weighted loss, a filter to significant interactions, a
precision-recall against called interactions, has to treat the 697{,}825 merged
doubles and 12{,}914 merged triples as carrying no p-value, and can recover the
sources' p-values only by joining back to the raw tables by gene pair, as this
document did.

\notesec{Changes the build would need}
Three, each at the step in Figure~\ref{fig:pipeline} that introduces the hazard.
The essentiality conversion should not produce a fitness that the deduplicator can
average with a measured one: either the converted 0 is kept as its own record with a
perturbation type the merge key distinguishes, or the conversion is dropped from
pools that also carry measured alleles of the gene. The deduplicator should keep the
source p-values beside the merged score, or combine them by a rule that reads them
(Fisher's method on the source p-values, or a normal test on the merged score with
the pooled source variance), rather than a $t$-test on two or three values. And a
merged record should carry which screens went into it, so that a within-screen
identity can be evaluated within a screen; the source tables carry the screen in the
strain identifier, and the build drops it at the aggregation step.
